\documentclass[12pt,a4paper]{amsart}

\usepackage{amsmath,amssymb,amsthm}
\usepackage{mathtools}
\usepackage[ngerman]{babel}
\usepackage{hyperref}
\usepackage{geometry}
\geometry{margin=2.5cm}

% -------------------------------------------------------
% Theorem-Umgebungen
% -------------------------------------------------------
\newtheorem{theorem}{Satz}[section]
\newtheorem{lemma}[theorem]{Lemma}
\newtheorem{corollary}[theorem]{Korollar}
\newtheorem{proposition}[theorem]{Proposition}
\theoremstyle{definition}
\newtheorem{definition}[theorem]{Definition}
\newtheorem{remark}[theorem]{Bemerkung}

% -------------------------------------------------------
% Eigene Befehle
% -------------------------------------------------------
\newcommand{\N}{\mathbb{N}}
\newcommand{\Z}{\mathbb{Z}}
\newcommand{\Q}{\mathbb{Q}}
\newcommand{\R}{\mathbb{R}}
\newcommand{\C}{\mathbb{C}}
\DeclareMathOperator{\Re}{Re}
\DeclareMathOperator{\Im}{Im}
\DeclareMathOperator{\res}{Res}

% -------------------------------------------------------
\begin{document}
% -------------------------------------------------------

\title{Die Riemann'sche Zeta-Funktion:\\
       Definition, analytische Fortsetzung und Funktionalgleichung}

\author{Michael Fuhrmann}
\address{Specialist Mathematics Project, Build 84}
\email{specialist-maths@localhost}
\date{\today}

\begin{abstract}
Die Riemann'sche Zeta-Funktion $\zeta(s)$ ist eines der tiefgr\"undigsten
Objekte in der gesamten Mathematik und verbindet Zahlentheorie, komplexe
Analysis und Physik.
Wir entwickeln die Theorie von Grund auf:
die Dirichlet-Reihen-Definition f\"ur $\Re(s) > 1$,
das Euler-Produkt \"uber Primzahlen,
die analytische Fortsetzung auf $\C \setminus \{1\}$
und Riemanns Funktionalgleichung
\[
  \zeta(s) \;=\; 2^s\,\pi^{s-1}\,\sin\!\Bigl(\tfrac{\pi s}{2}\Bigr)\,
  \Gamma(1-s)\,\zeta(1-s).
\]
Wir identifizieren die trivialen Nullstellen bei $s = -2, -4, -6, \ldots$
und f\"uhren die Riemann--Siegel-Formel $Z(t) = e^{i\theta(t)}\,\zeta\!\bigl(\tfrac{1}{2}+it\bigr)$
zur numerischen Nullstellensuche auf der kritischen Geraden ein.
Der kritische Streifen $0 < \Re(s) < 1$ ist die Arena, in der alle
nicht-trivialen Nullstellen liegen; die Riemann-Hypothese (noch offen)
behauptet, dass sie alle $\Re(s) = \tfrac{1}{2}$ erf\"ullen.
\end{abstract}

\maketitle
\tableofcontents

% ================================================================
\section{Einleitung}
% ================================================================

In seiner ber\"uhmten Abhandlung von 1859 \emph{\"Uber die Anzahl der
Primzahlen unter einer gegebenen Gr\"o\ss e} \cite{Riemann1859}
f\"uhrte Bernhard Riemann die komplexe Betrachtung der Funktion
\[
  \zeta(s) = \sum_{n=1}^{\infty} \frac{1}{n^s}
\]
ein, die Euler nur f\"ur reelles $s > 1$ untersucht hatte.
Riemanns entscheidende Einsicht war, $\zeta$ auf die ganze komplexe Ebene
fortzusetzen, ihre Funktionalgleichung herzuleiten und ihre Nullstellen mit
der Primzahlverteilung zu verbinden.

\begin{remark}
Wir verwenden die Standardnotation: $s = \sigma + it$ bezeichnet eine
komplexe Zahl mit Realteil $\sigma = \Re(s)$ und Imagin\"arteil $t = \Im(s)$.
Die Funktion $\Gamma$ bezeichnet Eulers Gamma-Funktion
(Abschnitt~\ref{sec:gamma}).
\end{remark}

% ================================================================
\section{Die Dirichlet-Reihe f\"ur \texorpdfstring{$\Re(s)>1$}{Re(s)>1}}
% ================================================================

\begin{definition}[Riemann'sche Zeta-Funktion -- Ausgangsbereich]
\label{def:zeta_series}
F\"ur $s \in \C$ mit $\Re(s) > 1$ ist die \emph{Riemann'sche Zeta-Funktion}
\[
  \zeta(s) \;=\; \sum_{n=1}^{\infty} \frac{1}{n^s}.
\]
\end{definition}

\begin{proposition}[Absolute Konvergenz]
\label{prop:convergence}
Die Reihe in Definition~\ref{def:zeta_series} konvergiert absolut und
gleichm\"a\ss ig auf jeder kompakten Teilmenge der Halbebene $\Re(s) > 1$.
Insbesondere ist $\zeta(s)$ holomorph auf $\{s : \Re(s) > 1\}$.
\end{proposition}

\begin{proof}
Sei $s = \sigma + it$ mit $\sigma > 1$. Dann gilt $|n^{-s}| = n^{-\sigma}$,
und $\sum_{n=1}^\infty n^{-\sigma}$ konvergiert mit dem Integraltest:
Das Integral $\int_1^\infty x^{-\sigma}\,dx = (\sigma-1)^{-1}$ ist endlich
f\"ur $\sigma > 1$.
Gleichm\"a\ss ige Konvergenz auf kompakten Mengen folgt aus dem
Weierstra\ss'schen $M$-Test mit $M_n = n^{-\sigma_0}$, wobei
$\sigma_0 = \min_K \Re(s) > 1$.
Holomorphie folgt aus der gleichm\"a\ss igen Konvergenz
holomorpher Funktionen.
\end{proof}

% ================================================================
\section{Das Euler-Produkt}
% ================================================================

\begin{theorem}[Euler-Produkt]
\label{thm:euler_product}
F\"ur $\Re(s) > 1$ gilt
\[
  \zeta(s) \;=\; \prod_{p\;\text{prim}} \frac{1}{1 - p^{-s}},
\]
wobei das Produkt absolut konvergiert.
\end{theorem}

\begin{proof}
Wir verwenden die geometrische Reihe: F\"ur jede Primzahl $p$ und $\Re(s)>1$:
\[
  \frac{1}{1 - p^{-s}} = \sum_{k=0}^{\infty} p^{-ks}.
\]
Das endliche Produkt \"uber Primzahlen $p \le P$ ergibt nach Ausmultiplizieren
-- unter Verwendung des Fundamentalsatzes der Arithmetik (eindeutige
Zerlegung) -- genau die Summe aller $n^{-s}$, deren Primfaktoren $\le P$ sind.
Der Fehler gegen\"uber $\zeta(s)$ ist durch $\sum_{n>P} n^{-\sigma}$
beschr\"ankt, was f\"ur $P \to \infty$ gegen $0$ strebt.
\end{proof}

\begin{corollary}
$\zeta(s) \ne 0$ f\"ur $\Re(s) > 1$.
\end{corollary}

\begin{proof}
Ein absolut konvergentes Produkt von Nicht-Null-Faktoren ist ungleich null.
\end{proof}

% ================================================================
\section{Die Gamma-Funktion}
\label{sec:gamma}
% ================================================================

\begin{definition}[Gamma-Funktion]
\label{def:gamma}
F\"ur $\Re(s) > 0$ ist Eulers \emph{Gamma-Funktion}
\[
  \Gamma(s) = \int_0^{\infty} t^{s-1} e^{-t}\,dt.
\]
\end{definition}

\begin{proposition}[Wichtige Eigenschaften von $\Gamma$]
\label{prop:gamma_props}
\begin{enumerate}
  \item[\textup{(i)}] \textit{Funktionalgleichung:} $\Gamma(s+1) = s\,\Gamma(s)$.
  \item[\textup{(ii)}] \textit{Fakult\"at:} $\Gamma(n) = (n-1)!$ f\"ur $n \in \N$.
  \item[\textup{(iii)}] \textit{Reflexionsformel:}
    $\Gamma(s)\,\Gamma(1-s) = \dfrac{\pi}{\sin(\pi s)}$
    f\"ur $s \notin \Z$.
  \item[\textup{(iv)}] \textit{Meromorphe Fortsetzung:} $\Gamma$ setzt sich zu
    einer meromorphen Funktion auf $\C$ fort mit einfachen Polen bei
    $s = 0, -1, -2, \ldots$
\end{enumerate}
\end{proposition}

\begin{proof}[Beweissskizze]
(i) Partielle Integration: $\int_0^\infty t^s e^{-t}\,dt
= [-t^s e^{-t}]_0^\infty + s\int_0^\infty t^{s-1}e^{-t}\,dt = s\,\Gamma(s)$.
(ii) Folgt aus (i) und $\Gamma(1) = 1$.
(iii) Klassische Identit\"at via Weierstra\ss'schem Produkt; siehe \cite{Ahlfors1979}.
(iv) Die Rekursion $\Gamma(s) = \Gamma(s+n)/[s(s+1)\cdots(s+n-1)]$ liefert
die meromorphe Fortsetzung.
\end{proof}

% ================================================================
\section{Analytische Fortsetzung von \texorpdfstring{$\zeta$}{zeta}}
% ================================================================

\begin{theorem}[Analytische Fortsetzung via Gamma-Integral]
\label{thm:continuation}
F\"ur $\Re(s) > 1$ gilt
\[
  \Gamma(s)\,\zeta(s) = \int_0^{\infty} \frac{t^{s-1}}{e^t - 1}\,dt.
\]
\end{theorem}

\begin{proof}
Substitution $t \mapsto nt$ in $\Gamma(s) = \int_0^\infty t^{s-1}e^{-t}\,dt$
liefert $\Gamma(s)/n^s = \int_0^\infty t^{s-1}e^{-nt}\,dt$.
Summation \"uber $n \ge 1$ (durch dominierte Konvergenz gerechtfertigt)
ergibt die Behauptung.
\end{proof}

\begin{theorem}[Einfacher Pol bei $s=1$]
\label{thm:pole}
$\zeta(s)$ hat einen einfachen Pol bei $s = 1$ mit Residuum $1$.
Genauer gilt:
\[
  \zeta(s) = \frac{1}{s-1} + \gamma + O(s-1),
\]
wobei $\gamma = 0{,}5772\ldots$ die Euler--Mascheroni-Konstante ist.
\end{theorem}

\begin{proof}[Beweisskizze]
Die Laurent-Entwicklung folgt aus der Euler--Maclaurin-Summenformel.
Der f\"uhrende Term liefert den einfachen Pol; der Konstantterm ergibt~$\gamma$.
\end{proof}

\begin{theorem}[Analytische Fortsetzung auf $\C \setminus \{1\}$]
\label{thm:full_continuation}
Die Funktion $\zeta(s)$ besitzt eine eindeutige analytische Fortsetzung
zu einer meromorphen Funktion auf ganz $\C$, mit einem einzigen einfachen
Pol bei $s = 1$ (Residuum~$1$) und keinen weiteren Singularit\"aten.
\end{theorem}

\begin{proof}[Beweisskizze]
Die Funktionalgleichung (Satz~\ref{thm:functional_eq}) dr\"uckt $\zeta(s)$
f\"ur $\Re(s) < 1$ durch $\zeta(1-s)$ mit $\Re(1-s) > 0$ aus und
liefert so die Fortsetzung in die linke Halbebene.
Eindeutigkeit folgt aus dem Identit\"atssatz f\"ur holomorphe Funktionen.
\end{proof}

% ================================================================
\section{Die Funktionalgleichung}
% ================================================================

\begin{theorem}[Riemann'sche Funktionalgleichung]
\label{thm:functional_eq}
Die meromorph fortgesetzte Funktion $\zeta$ erf\"ullt
\begin{equation}
  \label{eq:functional}
  \zeta(s) \;=\; 2^s\,\pi^{s-1}\,\sin\!\Bigl(\frac{\pi s}{2}\Bigr)\,
  \Gamma(1-s)\,\zeta(1-s)
  \qquad (s \in \C \setminus \{0,1\}).
\end{equation}
\"Aquivalent dazu: Definiert man
\[
  \xi(s) \;=\; \tfrac{1}{2}\,s(s-1)\,\pi^{-s/2}\,\Gamma\!\bigl(\tfrac{s}{2}\bigr)\,\zeta(s),
\]
so gilt $\xi(s) = \xi(1-s)$.
\end{theorem}

\begin{proof}[Beweisskizze]
Der Standardbeweis (nach Riemann 1859) verwendet die Jacobi-Theta-Funktion
\[
  \vartheta(t) = \sum_{n=-\infty}^{\infty} e^{-\pi n^2 t}
  = t^{-1/2}\,\vartheta(1/t) \quad (t > 0).
\]
Die Modularit\"atsidentit\"at $\vartheta(1/t) = t^{1/2}\vartheta(t)$ folgt
aus der Poisson-Summenformel.
Man leitet die vervollst\"andigte Zeta-Funktion
\[
  \xi(s) = \int_1^{\infty} \bigl(\vartheta(t)-1\bigr)
  \bigl(t^{s/2} + t^{(1-s)/2}\bigr)\frac{dt}{t}
\]
ab, die offensichtlich symmetrisch unter $s \leftrightarrow 1-s$ ist.
Ein vollst\"andiger Beweis findet sich in \cite{Davenport2000}, Kapitel~8.
\end{proof}

% ================================================================
\section{Triviale Nullstellen}
% ================================================================

\begin{theorem}[Triviale Nullstellen von $\zeta$]
\label{thm:trivial_zeros}
$\zeta(s) = 0$ gilt f\"ur $s = -2, -4, -6, \ldots$
(alle negativen geraden ganzen Zahlen).
Diese werden \emph{triviale Nullstellen} genannt.
\end{theorem}

\begin{proof}
Wir setzen $s = -2m$ f\"ur $m = 1, 2, 3, \ldots$ in die
Funktionalgleichung~\eqref{eq:functional} ein:
\[
  \zeta(-2m) = 2^{-2m}\,\pi^{-2m-1}\,\sin(-m\pi)\,\Gamma(2m+1)\,\zeta(2m+1).
\]
Da $\sin(-m\pi) = 0$ f\"ur $m \in \N$, folgt $\zeta(-2m) = 0$.
Da $\zeta(2m+1)$ f\"ur $2m+1 \ge 3$ endlich ist und $\Gamma(2m+1)$ keinen
Pol bei diesen Stellen hat, ist dies eine echte Nullstelle.

Die Einfachheit folgt daraus, dass $\sin(\pi s/2)$ eine einfache Nullstelle
bei $s = -2m$ hat und die anderen Faktoren dort weder Null noch unendlich sind.
\end{proof}

\begin{remark}
Die Werte bei negativen ungeraden ganzen Zahlen werden durch
Bernoulli-Zahlen $B_k$ gegeben:
\[
  \zeta(-n) = -\frac{B_{n+1}}{n+1} \quad (n = 1, 2, 3, \ldots).
\]
Insbesondere: $\zeta(-1) = -1/12$, $\zeta(-3) = 1/120$, $\zeta(0) = -1/2$.
\end{remark}

% ================================================================
\section{Der kritische Streifen und die Riemann-Hypothese}
% ================================================================

\begin{definition}[Kritischer Streifen und kritische Gerade]
\label{def:critical_strip}
Der \emph{kritische Streifen} ist $\{s \in \C : 0 < \Re(s) < 1\}$.
Die \emph{kritische Gerade} ist $\{s \in \C : \Re(s) = \tfrac{1}{2}\}$.

Die Nullstellen von $\zeta$ im kritischen Streifen hei\ss en
\emph{nicht-triviale Nullstellen}.
\end{definition}

\begin{proposition}[Nicht-triviale Nullstellen liegen im kritischen Streifen]
\label{prop:zeros_in_strip}
Alle nicht-trivialen Nullstellen von $\zeta(s)$ liegen im abgeschlossenen
Streifen $0 \le \Re(s) \le 1$.
\end{proposition}

\begin{proof}
F\"ur $\Re(s) > 1$: $\zeta(s) \ne 0$ nach dem Euler-Produkt.

F\"ur $\Re(s) < 0$ (triviale Nullstellen ausgenommen): Die Funktionalgleichung
zeigt, dass die einzigen Nullstellen von den Nullstellen von $\sin(\pi s/2)$
bei $s = -2, -4, \ldots$ stammen.
Insbesondere gilt $\zeta(0) = -1/2 \ne 0$, also ist $s=0$ keine Nullstelle.
\end{proof}

\begin{remark}[Die Riemann-Hypothese]
\label{rem:rh}
Riemann vermutete 1859, dass alle nicht-trivialen Nullstellen
$\Re(s) = \tfrac{1}{2}$ erf\"ullen.
Stand 2026 ist dies eines der sieben Millennium-Preisaufgaben
(Clay Mathematics Institute, 1.000.000~USD).
Mehr als $10^{13}$ Nullstellen wurden numerisch auf der kritischen
Geraden verifiziert \cite{Platt2021}.
\end{remark}

% ================================================================
\section{Die Riemann--Siegel-Formel}
% ================================================================

\begin{definition}[Riemann--Siegel-Theta-Funktion]
\label{def:theta}
F\"ur $t \in \R$ ist die \emph{Riemann--Siegel-Theta-Funktion}
\[
  \theta(t) \;=\; \arg\!\left[\Gamma\!\Bigl(\tfrac{1}{4} + \tfrac{it}{2}\Bigr)\right]
  - \frac{t}{2}\,\ln\pi.
\]
F\"ur gro\ss e $t$ gilt die asymptotische Entwicklung
\[
  \theta(t) = \frac{t}{2}\ln\frac{t}{2\pi} - \frac{t}{2} - \frac{\pi}{8}
  + \frac{1}{48t} - \frac{7}{5760\,t^3} + O(t^{-5}).
\]
\end{definition}

\begin{definition}[Hardy'sche $Z$-Funktion]
\label{def:Z_function}
Die \emph{Hardy'sche $Z$-Funktion} (Riemann--Siegel-$Z$-Funktion) ist
die reellwertige Funktion
\[
  Z(t) \;=\; e^{i\theta(t)}\,\zeta\!\Bigl(\tfrac{1}{2} + it\Bigr).
\]
\end{definition}

\begin{theorem}[Eigenschaften von $Z$]
\label{thm:Z_properties}
\begin{enumerate}
  \item[\textup{(i)}] $Z(t)$ ist reellwertig f\"ur alle $t \in \R$.
  \item[\textup{(ii)}] $|Z(t)| = |\zeta(\tfrac{1}{2} + it)|$.
  \item[\textup{(iii)}] Die Nullstellen von $Z(t)$ entsprechen genau
    den Nullstellen von $\zeta$ auf der kritischen Geraden.
  \item[\textup{(iv)}] Riemann--Siegel-Formel: F\"ur gro\ss es $t$:
    \[
      Z(t) = 2\sum_{n=1}^{N(t)} \frac{\cos(\theta(t) - t\ln n)}{\sqrt{n}}
      + R(t),
    \]
    wobei $N(t) = \lfloor\sqrt{t/(2\pi)}\rfloor$ und $R(t) = O(t^{-1/4})$.
\end{enumerate}
\end{theorem}

\begin{proof}
(i) Folgt aus der Funktionalgleichung:
$\overline{\zeta(\tfrac{1}{2}+it)} = \zeta(\tfrac{1}{2}-it)$.
Die Funktionalgleichung zeigt $e^{i\theta(t)}\zeta(\tfrac{1}{2}+it) \in \R$.

(ii) Wegen $|e^{i\theta(t)}| = 1$ gilt $|Z(t)| = |\zeta(\tfrac{1}{2}+it)|$.

(iii) Direkt aus (ii).

(iv) Die Riemann--Siegel-Formel entsteht aus der n\"aherungsweisen
Funktionalgleichung; vollst\"andige Herleitung in \cite{Titchmarsh1986},
Kapitel~4.
\end{proof}

% ================================================================
\section{Werte von \texorpdfstring{$\zeta$}{zeta} an positiven geraden Zahlen}
% ================================================================

\begin{theorem}[Eulers Formel f\"ur $\zeta(2k)$]
\label{thm:euler_values}
F\"ur $k = 1, 2, 3, \ldots$ gilt
\[
  \zeta(2k) \;=\; \frac{(-1)^{k+1}\,(2\pi)^{2k}\,B_{2k}}{2\,(2k)!},
\]
wobei $B_{2k}$ die Bernoulli-Zahlen sind. Insbesondere:
\[
  \zeta(2) = \frac{\pi^2}{6}, \quad
  \zeta(4) = \frac{\pi^4}{90}, \quad
  \zeta(6) = \frac{\pi^6}{945}.
\]
\end{theorem}

\begin{proof}[Beweisskizze]
Euler (1740) verwendete die Produktformel f\"ur $\sin$:
$\sin(\pi z) = \pi z \prod_{n=1}^{\infty}(1 - z^2/n^2)$.
Vergleich der Taylor-Entwicklung von $\ln\sin(\pi z)$ mit dem Logarithmus
des unendlichen Produkts liefert die Formel.
Vollst\"andiger Beweis in \cite{Apostol1976}.
\end{proof}

% ================================================================
\section{Schlussfolgerung}
% ================================================================

Wir haben die grundlegenden Eigenschaften der Riemann'schen Zeta-Funktion
entwickelt: Dirichlet-Reihe f\"ur $\Re(s)>1$, Euler-Produkt als Verbindung
zu den Primzahlen, analytische Fortsetzung auf $\C\setminus\{1\}$,
Funktionalgleichung~\eqref{eq:functional}, triviale Nullstellen bei negativen
geraden ganzen Zahlen und die Riemann--Siegel-Formel.

Die nicht-trivialen Nullstellen im kritischen Streifen halten den Schl\"ussel
zu pr\"azisen Absch\"atzungen f\"ur die Primzahlverteilung bereit
(via expliziter Formel; vgl.\ Paper~23).
Ob alle nicht-trivialen Nullstellen auf der kritischen Geraden $\Re(s) = 1/2$
liegen -- die Riemann-Hypothese -- bleibt das ber\"uhmteste offene Problem
der Mathematik.

% ================================================================
\begin{thebibliography}{10}

\bibitem{Riemann1859}
B.~Riemann,
\emph{\"Uber die Anzahl der Primzahlen unter einer gegebenen Gr\"o\ss e},
Monatsberichte der Berliner Akademie (1859), 671--680.
Neu abgedruckt in: \emph{Gesammelte Werke}, Teubner, Leipzig, 1892.

\bibitem{Davenport2000}
H.~Davenport,
\emph{Multiplicative Number Theory}, 3.~Aufl.\
(revidiert von H.~L.~Montgomery),
Graduate Texts in Mathematics~74,
Springer, New York, 2000.

\bibitem{Titchmarsh1986}
E.~C.~Titchmarsh,
\emph{The Theory of the Riemann Zeta-Function}, 2.~Aufl.\
(revidiert von D.~R.~Heath-Brown),
Oxford University Press, 1986.

\bibitem{Ahlfors1979}
L.~V.~Ahlfors,
\emph{Complex Analysis}, 3.~Aufl.,
McGraw-Hill, New York, 1979.

\bibitem{Apostol1976}
T.~M.~Apostol,
\emph{Introduction to Analytic Number Theory},
Undergraduate Texts in Mathematics,
Springer, New York, 1976.

\bibitem{Platt2021}
D.~J.~Platt und T.~S.~Trudgian,
The Riemann hypothesis is true up to $3 \times 10^{12}$,
\emph{Bull.\ London Math.\ Soc.}\ \textbf{53} (2021), 792--797.

\end{thebibliography}

\end{document}
